\Article{29}
Es bezeichne \(t\) eine gegebene lineare Funktion der Größen \(x,y,z\) etc.; es sei etwa
\[
t=fx+gy+hz+\text{etc.}+k
\]
Der Wert von \(t\), der sich aus den plausibelsten Werten von \(x,y,z\) etc. ergibt, ist dann \(=fA+gB+hC+\text{etc.}+k\); wir bezeichnen ihn mit \(K\). Wird dieser als wahrer Wert von \(t\) angenommen, so entsteht ein Fehler, für den sich
\SourceLeaf{040}
\[
=fEx+gEy+hEz+\text{etc.}
\]
ergibt und den wir mit \(Et\) bezeichnen. Offenbar ist der mittlere Wert dieses Fehlers \(=0\), oder der Fehler ist frei von einem konstanten Anteil. Der mittlere Wert des Quadrats \((Et)^2\) oder der mittlere Wert der Summe
\[
\begin{array}{llll}
ff(Ex)^2&+2fg\,Ex.Ey&+2fh\,Ex.Ez&+\text{etc.}\\
&+gg(Ey)^2&+2gh\,Ey.Ez&+\text{etc.}\\
&&+hh(Ez)^2&+\text{etc. etc.}
\end{array}
\]
ist aber nach dem im vorigen Artikel Dargelegten gleich dem Produkt aus \(mmp\) und der Summe
\[
\begin{array}{llll}
ff[\alpha\alpha]&+2fg[\alpha\beta]&+2fh[\alpha\gamma]&+\text{etc.}\\
&+gg[\beta\beta]&+2gh[\beta\gamma]&+\text{etc.}\\
&&+hh[\gamma\gamma]&+\text{etc. etc.}
\end{array}
\]
oder dem Produkt aus \(mmp\) und dem Wert der Funktion \(\Omega-M\), der sich durch die Einsetzungen
\[
\xi=f,\quad\eta=g,\quad\zeta=h\quad\text{etc.}
\]
ergibt. Bezeichnen wir also diesen bestimmten Wert der Funktion \(\Omega-M\) mit \(\omega\), so ist der zu befürchtende mittlere Fehler, solange wir an der Bestimmung \(t=K\) festhalten, \(=m\BareRoot{p\omega}\) oder das Gewicht dieser Bestimmung ist \(=\frac1\omega\).

Da identisch
\[
\Omega-M=(x-A)\xi+(y-B)\eta+(z-C)\zeta+\text{etc.}
\]
gilt, ist klar, dass \(\omega\) auch gleich dem bestimmten Wert des Ausdrucks
\[
(x-A)f+(y-B)g+(z-C)h+\text{etc.}
\]
oder dem bestimmten Wert von \(t-K\) ist, der sich ergibt, wenn den Unbestimmten \(x,y,z\) etc. diejenigen Werte zugewiesen werden, die den Werten \(\xi,\eta,\zeta\) etc. gleich \(f,g,h\) etc. entsprechen.

Schließlich bemerken wir: Wird \(t\) identisch als Funktion von \(\xi,\eta,\zeta\) etc. dargestellt, so ist sein konstanter Anteil notwendig \(=K\). Gilt also identisch
\[
t=F\xi+G\eta+H\zeta+\text{etc.}+K,\quad\text{so gilt}\quad\omega=fF+gG+hH+\text{etc.}
\]
\SourceLeaf{041}
\Article{30}
Die Funktion \(\Omega\) erreicht, wie wir oben gesehen haben, ihren \textit{absoluten Minimalwert} \(M\), wenn wir \(x=A,\allowbreak y=B,\allowbreak z=C\) etc. oder \(\xi=0,\allowbreak\eta=0,\allowbreak\zeta=0\) etc. setzen. Ist jedoch einer dieser Größen bereits ein \textit{anderer} Wert zugewiesen, z. B. \(x=A+\Delta\), so kann \(\Omega\) bei Veränderung der übrigen Größen einen \textit{relativen Minimalwert} erreichen, der sich offenbar mithilfe der Gleichungen
\[
x=A+\Delta,\quad\frac{d\Omega}{dy}=0,\quad\frac{d\Omega}{dz}=0\quad\text{etc.}
\]
ergibt. Es muss also \(\eta=0,\zeta=0\) etc. gelten, und aus der Beziehung ergibt sich folglich
\[
x=A+[\alpha\alpha]\xi+[\alpha\beta]\eta+[\alpha\gamma]\zeta+\text{etc.},\quad\xi=\frac{\Delta}{[\alpha\alpha]}
\]
Zugleich haben wir
\[
y=B+\frac{[\alpha\beta]\Delta}{[\alpha\alpha]},\quad z=C+\frac{[\alpha\gamma]\Delta}{[\alpha\alpha]}\quad\text{etc.}
\]
Der relative Minimalwert von \(\Omega\) ist dann \(=[\alpha\alpha]\xi\xi+M=M+\frac{\Delta\Delta}{[\alpha\alpha]}\). Umgekehrt schließen wir daraus: Darf der Wert von \(\Omega\) die vorgegebene Grenze \(M+\mu\mu\) nicht überschreiten, so muss der Wert von \(x\) notwendig zwischen den Grenzen \(A-\mu\BareRoot{[\alpha\alpha]}\) und \(A+\mu\BareRoot{[\alpha\alpha]}\) liegen. Es verdient Beachtung, dass \(\mu\BareRoot{[\alpha\alpha]}\) gleich dem zu befürchtenden mittleren Fehler im plausibelsten Wert von \(x\) ist, wenn wir \(\mu=m\BareRoot{p}\) setzen, d. h., wenn \(\mu\) gleich dem mittleren Fehler solcher Beobachtungen ist, denen das Gewicht \(=1\) zugewiesen wird.

Untersuchen wir allgemeiner den Minimalwert von \(\Omega\), der bei einem gegebenen Wert von \(t\) auftreten kann, wobei \(t\) wie im vorigen Artikel die lineare Funktion \(fx+gy+hz+\text{etc.}+k\) bezeichnet und ihr plausibelster Wert \(=K\) ist: Der vorgegebene Wert von \(t\) werde mit \(K+\kappa\) bezeichnet. Aus der Theorie der Maxima und Minima ist bekannt, dass die Lösung des Problems aus den Gleichungen
\[
\begin{aligned}
\frac{d\Omega}{dx}&=\theta\frac{dt}{dx}\\[3pt]
\frac{d\Omega}{dy}&=\theta\frac{dt}{dy}\\[3pt]
\frac{d\Omega}{dz}&=\theta\frac{dt}{dz}\quad\text{etc.}
\end{aligned}
\]
zu suchen ist, oder \(\xi=\theta f,\allowbreak\eta=\theta g,\allowbreak\zeta=\theta h\) etc., wobei \(\theta\) einen noch unbestimmten Multiplikator bezeichnet. Setzen wir daher wie im vorigen Artikel \textit{identisch}
\[
t=F\xi+G\eta+H\zeta+\text{etc.}+K
\]
so haben wir
\SourceLeaf{042}
\[
\begin{aligned}
K+\kappa&=\theta(fF+gG+hH+\text{etc.})+K,\quad\text{oder}\\[3pt]
\theta&=\frac{\kappa}{\omega}
\end{aligned}
\]
wobei \(\omega\) dieselbe Bedeutung wie im vorigen Artikel hat. Da ferner \(\Omega-M\) identisch eine homogene Funktion zweiten Grades der Unbestimmten \(\xi,\eta,\zeta\) etc. ist, ergibt sich von selbst, dass ihr Wert für \(\xi=\theta f,\allowbreak\eta=\theta g,\allowbreak\zeta=\theta h\) etc. zu \(=\theta\theta\omega\) wird und folglich der Minimalwert, den \(\Omega\) für \(t=K+\kappa\) erreichen kann, zu \(=M+\theta\theta\omega=M+\frac{\kappa\kappa}{\omega}\). Umgekehrt muss, wenn \(\Omega\) einen vorgegebenen Wert \(M+\mu\mu\) nicht überschreiten darf, der Wert von \(t\) notwendig zwischen den Grenzen \(K-\mu\BareRoot{\omega}\) und \(K+\mu\BareRoot{\omega}\) liegen; dabei ist \(\mu\BareRoot{\omega}\) gleich dem zu befürchtenden mittleren Fehler bei der plausibelsten Bestimmung von \(t\), wenn für \(\mu\) der mittlere Fehler der Beobachtungen genommen wird, denen das Gewicht \(=1\) zugewiesen wird.

\clearpage
